\documentclass[12pt,a4paper]{article}

\usepackage[utf8]{inputenc}
\usepackage[T1]{fontenc}
\usepackage[french]{babel}
\usepackage{geometry}
\usepackage{booktabs}
\usepackage{siunitx}
\usepackage{array}
\usepackage{xcolor}
\usepackage{titlesec}
\usepackage{fancyhdr}
\usepackage{pgfplots}

\geometry{margin=2.3cm}
\definecolor{bleufonce}{HTML}{123B63}
\definecolor{bleuclair}{HTML}{2F80C1}
\definecolor{turquoise}{HTML}{159A9C}
\definecolor{orange}{HTML}{E18A36}
\definecolor{violet}{HTML}{7851A9}
\definecolor{grisclair}{HTML}{F1F5F9}
\pgfplotsset{compat=1.17}
\titleformat{\section}
  {\Large\bfseries\color{bleufonce}}
  {\thesection}{0.7em}{}
  [\vspace{0.2em}{\color{bleuclair}\titlerule}]
\titleformat{\subsection}
  {\large\bfseries\color{bleuclair}}
  {\thesubsection}{0.7em}{}
\pagestyle{fancy}
\fancyhf{}
\fancyhead[L]{\small\color{bleufonce}Benchmark du tas max}
\fancyhead[R]{\small\color{bleufonce}ERRAMI Youssef}
\fancyfoot[C]{\color{bleufonce}\thepage}
\setlength{\headheight}{15pt}
\sisetup{
  output-decimal-marker = {,},
  group-separator = {\,},
  group-minimum-digits = 4
}

\title{Comparaison de conteneurs pour l'implémentation d'un tas max}
\author{ERRAMI Youssef}
\date{30 octobre 2026}

\begin{document}
\begin{center}
  {\color{bleuclair}\rule{\textwidth}{1.2pt}}\\[0.8em]
  {\Huge\bfseries\color{bleufonce}
    Comparaison de conteneurs\\[0.2em]
    pour un tas max}\\[0.6em]
  {\Large\color{bleuclair}Comparaison des performances en C++}\\[1.2em]
  {\large\bfseries ERRAMI Youssef}\\[0.3em]
  {\color{bleufonce}30 octobre 2026}\\[0.8em]
  {\color{bleuclair}\rule{\textwidth}{1.2pt}}
\end{center}
\vspace{0.8em}

\noindent
\colorbox{grisclair}{%
  \parbox{\dimexpr\textwidth-2\fboxsep\relax}{%
    \textcolor{bleufonce}{\textbf{Comparaison.}}
    Tas max maison sur \texttt{vector}, \texttt{list} et \texttt{map},
    comparés au tas max STL sur \texttt{vector}.%
  }%
}

\section{Ordinateur}
\begin{itemize}
  \item MSI Stealth GS66 12UGS.
  \item Intel Core i7-12700H, 14 coeurs.
  \item NVIDIA GeForce RTX 3070 Ti.
  \item 16 Go de mémoire DDR5.
\end{itemize}

\section{Résultats}

Les temps sont exprimés en millisecondes. Les résultats pour $N\leq 5\,000$
et $N=1\,000\,000$ sont les médianes de cinq exécutions. À
$N=100\,000\,000$, une seule exécution est disponible : les chiffres ne sont
donc pas des médianes. L'extraction mesure le retrait successif des $N$
maximums jusqu'à vider le tas. Les mêmes valeurs aléatoires sont utilisées
par les variantes d'un même essai. « -- » signifie qu'aucune mesure n'est
disponible. Les variantes maison sur \texttt{list} et \texttt{map} sont
mesurées seulement jusqu'à $N=5\,000$.

\subsection{Insertion}

\begin{table}[htbp]
\centering
\caption{Tas max : temps d'insertion (ms)}
\label{tab:max-insertion}
\small
\begin{tabular}{r S[table-format=6.6] S[table-format=6.6] S[table-format=6.6] S[table-format=6.6]}
\toprule
{$N$} & {\texttt{maison/vector}} & {\texttt{maison/list}} & {\texttt{maison/map}} & {\texttt{STL/vector}} \\
\midrule
100       & 0.001500 & 0.018400 & 0.023200 & 0.001600 \\
500       & 0.007100 & 0.508900 & 0.133100 & 0.006800 \\
1000      & 0.013000 & 1.986900 & 0.312800 & 0.012700 \\
2000      & 0.025400 & 9.498000 & 0.589800 & 0.025500 \\
5000      & 0.063200 & 98.962800 & 1.483000 & 0.060200 \\
10000     & \multicolumn{1}{c}{--} & \multicolumn{1}{c}{--} & \multicolumn{1}{c}{--} & \multicolumn{1}{c}{--} \\
100000    & \multicolumn{1}{c}{--} & \multicolumn{1}{c}{--} & \multicolumn{1}{c}{--} & \multicolumn{1}{c}{--} \\
1000000   & 12.620200 & \multicolumn{1}{c}{--} & \multicolumn{1}{c}{--} & 12.064200 \\
10000000  & \multicolumn{1}{c}{--} & \multicolumn{1}{c}{--} & \multicolumn{1}{c}{--} & \multicolumn{1}{c}{--} \\
100000000 & 2252.815600 & \multicolumn{1}{c}{--} & \multicolumn{1}{c}{--} & 2210.443300 \\
\bottomrule
\end{tabular}
\end{table}

\subsection{Extraction du maximum}

\begin{table}[htbp]
\centering
\caption{Tas max : temps d'extraction des $N$ maximums (ms)}
\label{tab:max-extraction}
\small
\begin{tabular}{r S[table-format=6.6] S[table-format=6.6] S[table-format=6.6] S[table-format=6.6]}
\toprule
{$N$} & {\texttt{maison/vector}} & {\texttt{maison/list}} & {\texttt{maison/map}} & {\texttt{STL/vector}} \\
\midrule
100       & 0.003900 & 0.026700 & 0.041200 & 0.002400 \\
500       & 0.026800 & 0.890000 & 0.313400 & 0.013500 \\
1000      & 0.052500 & 4.029200 & 0.908500 & 0.022800 \\
2000      & 0.113700 & 18.773400 & 2.196400 & 0.056500 \\
5000      & 0.308800 & 170.574900 & 7.886300 & 0.125600 \\
10000     & \multicolumn{1}{c}{--} & \multicolumn{1}{c}{--} & \multicolumn{1}{c}{--} & \multicolumn{1}{c}{--} \\
100000    & \multicolumn{1}{c}{--} & \multicolumn{1}{c}{--} & \multicolumn{1}{c}{--} & \multicolumn{1}{c}{--} \\
1000000   & 161.764400 & \multicolumn{1}{c}{--} & \multicolumn{1}{c}{--} & 158.403400 \\
10000000  & \multicolumn{1}{c}{--} & \multicolumn{1}{c}{--} & \multicolumn{1}{c}{--} & \multicolumn{1}{c}{--} \\
100000000 & 74721.719500 & \multicolumn{1}{c}{--} & \multicolumn{1}{c}{--} & 182454.176300 \\
\bottomrule
\end{tabular}
\end{table}

\clearpage
\section{Graphiques}

Les courbes relient les temps mesurés pour $N=100$, $500$, $1\,000$,
$2\,000$, $5\,000$, $1\,000\,000$ et $100\,000\,000$. Les segments entre
points servent uniquement à visualiser l'évolution : ils ne représentent pas
de mesures supplémentaires. L'axe vertical logarithmique rend visibles les
temps courts et longs.

\pgfplotsset{
  benchmark axis/.style={
    width=0.92\textwidth,
    height=7.0cm,
    xmode=linear,
    ymode=log,
    log basis y=10,
    ymin=0.001,
    ymax=300000,
    xmin=0.7,
    xmax=10.3,
    xtick={1,2,3,4,5,6,7,8,9,10},
    xticklabels={100,500,1k,2k,5k,10k,100k,1M,10M,100M},
    xlabel={Nombre d'éléments $N$},
    ylabel={Temps mesuré (ms, échelle logarithmique)},
    ytick={0.001,0.01,0.1,1,10,100,1000,10000,100000},
    yticklabels={$10^{-3}$,$10^{-2}$,$10^{-1}$,$10^0$,$10^1$,$10^2$,$10^3$,$10^4$,$10^5$},
    scaled y ticks=false,
    grid=both,
    major grid style={gray!30},
    minor grid style={gray!12},
    tick label style={font=\small},
    label style={font=\small\color{bleufonce}},
    legend style={
      at={(0.5,-0.20)},
      anchor=north,
      legend columns=2,
      draw=none,
      font=\small
    },
    line width=1.1pt,
    mark size=3.0pt
  }
}

\subsection{Insertion}
\begin{figure}[htbp]
\centering
\begin{tikzpicture}
\begin{axis}[benchmark axis]
\addplot[color=bleufonce,mark=*] coordinates {
  (1,0.0015) (2,0.0071) (3,0.0130) (4,0.0254) (5,0.0632)
  (8,12.6202) (10,2252.8156)
};
\addlegendentry{maison/vector}
\addplot[color=turquoise,mark=square*] coordinates {
  (1,0.0184) (2,0.5089) (3,1.9869) (4,9.4980) (5,98.9628)
};
\addlegendentry{maison/list}
\addplot[color=orange,mark=triangle*] coordinates {
  (1,0.0232) (2,0.1331) (3,0.3128) (4,0.5898) (5,1.4830)
};
\addlegendentry{maison/map}
\addplot[color=violet,mark=diamond*] coordinates {
  (1,0.0016) (2,0.0068) (3,0.0127) (4,0.0255) (5,0.0602)
  (8,12.0642) (10,2210.4433)
};
\addlegendentry{STL/vector}
\end{axis}
\end{tikzpicture}
\caption{Insertion dans les tas max.}
\label{fig:max-insertion}
\end{figure}

\subsection{Extraction}
\begin{figure}[htbp]
\centering
\begin{tikzpicture}
\begin{axis}[benchmark axis]
\addplot[color=bleufonce,mark=*] coordinates {
  (1,0.0039) (2,0.0268) (3,0.0525) (4,0.1137) (5,0.3088)
  (8,161.7644) (10,74721.7195)
};
\addlegendentry{maison/vector}
\addplot[color=turquoise,mark=square*] coordinates {
  (1,0.0267) (2,0.8900) (3,4.0292) (4,18.7734) (5,170.5749)
};
\addlegendentry{maison/list}
\addplot[color=orange,mark=triangle*] coordinates {
  (1,0.0412) (2,0.3134) (3,0.9085) (4,2.1964) (5,7.8863)
};
\addlegendentry{maison/map}
\addplot[color=violet,mark=diamond*] coordinates {
  (1,0.0024) (2,0.0135) (3,0.0228) (4,0.0565) (5,0.1256)
  (8,158.4034) (10,182454.1763)
};
\addlegendentry{STL/vector}
\end{axis}
\end{tikzpicture}
\caption{Extraction des maximums jusqu'à vider le tas max.}
\label{fig:max-extraction}
\end{figure}

\section{Comparaison}

À $N=1\,000\,000$, les médianes d'insertion sont de \SI{12.6202}{ms} pour
la version maison/vector et \SI{12.0642}{ms} pour la STL/vector. L'extraction
complète prend respectivement \SI{161.7644}{ms} et \SI{158.4034}{ms}.

À $N=100\,000\,000$, une seule exécution a été conservée : insertion
\SI{2252.8156}{ms} (maison/vector) contre \SI{2210.4433}{ms} (STL/vector) ;
extraction \SI{74721.7195}{ms} contre \SI{182454.1763}{ms}. L'écart important
observé à l'extraction doit être interprété avec prudence, car ce résultat
n'est pas confirmé par des répétitions. Les mesures de ce rapport sont
disponibles dans \texttt{benchmark\_resultats.csv}.

\end{document}
